\chapter{Random Walks}

\begin{center}\itshape
On the line, a symmetric walk returns to the origin infinitely often.\\
In dimension three or more, it leaves and the expected number of returns is finite.
\end{center}

Let $X_1,X_2,\ldots$ be independent fair signs, $P(X_i=\pm1)=1/2$, built on disjoint bits. The simple random walk is $S_n=X_1+\cdots+X_n$, with $S_0=0$. It is periodic of period $2$: after an even number of steps the position is even. Theorem~\ref{thm:local} gives $P(S_{2m}=0)\sim 1/\sqrt{\pi m}$.

\section{Reflection and the maximum}

Let $M_n=\max_{0\le k\le n}S_k$.

\begin{theorem}[reflection]\label{thm:rwreflect}
Let $a$ be a positive integer and let $b$ be an integer with the same parity as $n$. The number of paths of length $n$ from $0$ to $b$ that visit $a$ at least once equals the number of paths of length $n$ from $0$ to $2a-b$, provided one counts a path that ends at or above $a$ as having visited $a$. Equivalently, for every $b\le a$,
\[
 P(M_n\ge a,\,S_n=b)=P(S_n=2a-b).
\]
\end{theorem}
\begin{proof}
Let $\tau$ be the first time a path visits $a$. Reflect the positions on $\{0,\ldots,\tau\}$ across the level $a$, and leave the rest of the path unchanged. The new path runs from $2a$ to $b$. Subtract $2a$ from every position: the steps are unchanged, and the path now runs from $0$ to $b-2a$. Flip every sign. The path runs from $0$ to $2a-b$, and the construction is reversible on the set of paths from $0$ to $2a-b$. When $b\le a$ one has $2a-b\ge a$, so every path from $0$ to $2a-b$ does visit $a$, and the inverse is defined. The two sets of paths are therefore in bijection. Divide by $2^n$.
\end{proof}

Summing over $b\le a$ of the same parity,
\[
 P(M_n\ge a)=P(S_n\ge a)+P(S_n>a),
\]
because $\{S_n\ge a\}$ is already contained in $\{M_n\ge a\}$, and the reflected copies account for the paths that touch $a$ and finish strictly below it, which are in bijection with $\{S_n>a\}$.

\section{Gambler's ruin}

Two players share $N$ units. The walk starts at $k\in\{0,\ldots,N\}$ and stops at $0$ or at $N$. Let $h_k$ be the probability of hitting $N$ before $0$.

\begin{theorem}[ruin]\label{thm:ruin}
For a fair walk, $h_k=k/N$. If each step is $+1$ with probability $p\in(0,1)\setminus\{1/2\}$ and $-1$ with probability $q=1-p$, then
\[
 h_k=\frac{(q/p)^k-1}{(q/p)^N-1}.
\]
\end{theorem}
\begin{proof}
Condition on the first step. Then $h_0=0$, $h_N=1$, and $h_k=p h_{k+1}+q h_{k-1}$ for $0<k<N$. For $p=1/2$ the second difference vanishes, so $h_k$ is linear, and the boundary values give $h_k=k/N$. For $p\ne1/2$ the characteristic equation $p r^2-r+q=0$ has roots $1$ and $q/p$. The general solution is $A+B(q/p)^k$. The boundary values fix $A$ and $B$.
\end{proof}

The same conditioning gives the expected duration. For the fair walk, if $e_k$ is the expected number of steps to absorption from $k$, then $e_0=e_N=0$ and $e_k=1+(e_{k+1}+e_{k-1})/2$. The solution is $e_k=k(N-k)$. Check: the second difference of $k(N-k)$ is $-2$, so adding $1$ in the equation matches.

\section{Recurrence and transience}

Let $u_n=P(S_n=0)$ in $\Z^d$, where the walk steps to one of the $2d$ neighbours with equal probability $1/(2d)$. Let $f$ be the probability of a return to $0$ at some positive time.

\begin{proposition}[expected number of visits]\label{prop:visits}
Counting the visit at time $0$, the expected number of visits to the origin is $\sum_{n\ge0}u_n$, and this sum equals $1/(1-f)$ when $f<1$. It is infinite when $f=1$.
\end{proposition}
\begin{proof}
After each return the walk starts afresh, because the increments after the return are a copy of the original sequence, independent of the past. The number of visits is therefore geometric: it equals $m$ with probability $f^{m-1}(1-f)$ when $f<1$, including the visit at time $0$. The expectation is $1/(1-f)$. If $f=1$ the walk returns almost surely, and then returns again, so there are infinitely many visits and the expectation is infinite. In every case the expectation equals $\sum_n u_n$ by writing the number of visits as $\sum_n\mathbf 1_{S_n=0}$ and applying monotone convergence of nonnegative expectations, which is the tail-sum idea of Volume~1 applied to the partial sums.
\end{proof}

\begin{theorem}[P\'olya]\label{thm:polya}
The simple random walk on $\Z^d$ is recurrent for $d=1$ and $d=2$, and transient for every $d\ge3$.
\end{theorem}
\begin{proof}
In dimension $1$, $u_{2m}\sim 1/\sqrt{\pi m}$ by Theorem~\ref{thm:local}, and the odd terms vanish. The series $\sum m^{-1/2}$ diverges, so $\sum u_n=\infty$ and $f=1$.

In dimension $2$, a return to the origin in $2m$ steps requires $m$ horizontal steps and $m$ vertical steps, and within the horizontal steps equally many left and right. Counting gives
\[
 u_{2m}=\Bigl(\binom{2m}{m}2^{-2m}\Bigr)^2.
\]
Theorem~\ref{thm:local} turns this into $\sim 1/(\pi m)$. The harmonic series diverges, so the walk is recurrent.

In dimension $d\ge1$ the inversion formula on the torus reads
\[
 u_n=(2\pi)^{-d}\int_{[-\pi,\pi]^d}\Bigl(\frac1d\sum_{j=1}^d\cos\theta_j\Bigr)^n\dd\theta.
\]
It holds because $\int_{-\pi}^{\pi}e^{im\varphi}\,\dd\varphi$ equals $2\pi$ if $m=0$ and equals $0$ otherwise, and expanding the $n$th power enumerates the walks. Write $\phi(\theta)=d^{-1}\sum_j\cos\theta_j$. Then $|\phi(\theta)|\le1$, with equality in $[-\pi,\pi]^d$ only at the origin. On the complement of a neighbourhood of $0$, $|\phi|\le\rho<1$, so that contribution to $u_n$ is $O(\rho^n)$. Near the origin, $\cos u=1-u^2/2+O(u^4)$, so $\phi(\theta)=1-|\theta|^2/(2d)+O(|\theta|^4)$. For small $\theta$ and large $n$, $\phi(\theta)^n\le\exp(-c n|\theta|^2)$. The integral of that Gaussian over $\R^d$ is of order $n^{-d/2}$. Thus $u_n=O(n^{-d/2})$. For $d\ge3$ the series $\sum n^{-d/2}$ converges, so $\sum u_n<\infty$ and $f<1$.
\end{proof}

The one-dimensional walk returns, but the expected time of the first return is infinite: the probability of a first return at time $2m$ is asymptotic to $1/\sqrt{\pi}\,m^{-3/2}$, by the reflection count $f_{2m}=u_{2m-2}/(2m-1)$, and $\sum m\cdot m^{-3/2}$ diverges. The same $m^{-3/2}$ is why the duration of the fair game in Theorem~\ref{thm:ruin}, although finite for fixed $N$, grows like $N^2$.

\section{Harmonic functions}

On a finite connected graph, let the walk step to a uniform neighbour. For two vertices $a\ne b$, let $h(x)$ be the probability of hitting $b$ before $a$, starting from $x$.

\begin{proposition}[harmonic hitting probabilities]\label{prop:harmonic}
The function $h$ is the unique function with $h(a)=0$, $h(b)=1$, and $h(x)$ equal to the average of $h$ over the neighbours of $x$, for every $x\notin\{a,b\}$.
\end{proposition}
\begin{proof}
The boundary values are the definition of absorption. At any other $x$, the first step goes to a uniform neighbour and the Markov property of independent steps gives the average. If $h_1$ and $h_2$ both solve the problem, their difference $d$ is zero at $a$ and at $b$ and is harmonic elsewhere. A harmonic function on a finite connected graph attains its maximum. If the maximum were positive it could not occur at $a$ or $b$, and at an interior maximum the average of the neighbours would be strictly smaller unless $d$ is constant, which it cannot be if it is zero at $a$. So $d\le0$. The same argument for $-d$ gives $d\ge0$.
\end{proof}

Exercise~\ref{ex:electric} reads $h$ as a voltage.

The walk on the line is the fair case of Theorem~\ref{thm:ruin}. The next chapter allows the step law to depend on the current state, as long as the state space is finite.

\section{Exercises}

\begin{exercise}[$\star$]\label{ex:ballotwalk}
Deduce Theorem~\ref{thm:ballot} from Theorem~\ref{thm:rwreflect} by counting paths that stay positive, or explain in the other direction which hypothesis of the reflection theorem is the ballot hypothesis $a>b$.
\end{exercise}

\begin{exercise}[$\star\star$]\label{ex:maxlaw}
For the fair walk on $\Z$, show that $P(M_n\ge a,\,S_n\le a)=P(S_n\ge a)$ is not the claim. Using Theorem~\ref{thm:rwreflect}, show $P(M_n\ge a)=P(S_n\ge a)+P(S_n>a)$.
\end{exercise}

\begin{exercise}[$\star\star$]\label{ex:timepositive}
Show that the probability of a first return to $0$ at time $2m$ equals $u_{2m}/(2m-1)$, for $m\ge1$, with $u_{2m}=P(S_{2m}=0)$.
\end{exercise}

\begin{exercise}[$\star$]\label{ex:unfair}
In Theorem~\ref{thm:ruin}, with $p=2/3$ and $N=3$, compute $h_1$ and $h_2$ from the closed form and from the linear system.
\end{exercise}

\begin{exercise}[$\star\star$]\label{ex:meet}
Two independent fair walks on $\Z$ start at $0$ and at $2$. Show that they meet almost surely, by looking at their difference.
\end{exercise}

\begin{exercise}[$\star\star$]\label{ex:escape}
Show that $\sum_{m\ge1}m\,P(\text{first return at }2m)=\infty$ for the fair walk on $\Z$.
\end{exercise}

\begin{exercise}[$\star\star$]\label{ex:exittime}
Let $\tau$ be the first time $|S_n|\ge R$ for the fair walk on $\Z$, with $R$ a positive integer. Show that $\mathbb E[\tau]=R^2$.
\end{exercise}

\begin{exercise}[$\star\star$]\label{ex:electric}
On a finite connected graph, let $h(x)$ be the probability that the simple random walk from $x$ hits $b$ before $a$.
\begin{enumerate}[label=(\alph*)]
\item Recall that $h$ is the unique harmonic function with $h(a)=0$ and $h(b)=1$.
\item Let the edges be unit resistors. Explain why $h$ is the voltage when $a$ is earthed and $b$ is held at one volt: current into a vertex other than $a$ and $b$ is zero, and current along an edge is the difference of voltages.
\item Let $R$ be the effective resistance between $a$ and $b$, the reciprocal of the total current leaving $b$. Show that the walk from $a$ hits $b$ before returning to $a$ with probability $1/(\deg(a)R)$.
\item Compute $R$ for a path of length $m$, and for a cycle in which $a$ and $b$ are separated by paths of lengths $m$ and $k$.
\end{enumerate}
\end{exercise}

\section{Hints}

\begin{hint}{ex:ballotwalk}
A path that stays strictly positive is a path whose maximum condition is tied to the end.
\end{hint}
\begin{hint}{ex:maxlaw}
Sum $P(S_n=2a-b)$ over $b\le a$.
\end{hint}
\begin{hint}{ex:timepositive}
A first return at $2m$ is a path from $0$ to $0$ that stays positive in between, up to a factor $2$ for the sign. The reflection count gives $f_{2m}=u_{2m-2}-u_{2m}$.
\end{hint}
\begin{hint}{ex:unfair}
$(q/p)=1/2$.
\end{hint}
\begin{hint}{ex:meet}
The difference of two fair walks, observed at joint steps, is a walk that can stay put or move by $2$. It still visits every site.
\end{hint}
\begin{hint}{ex:escape}
Use $f_{2m}\sim c\,m^{-3/2}$.
\end{hint}
\begin{hint}{ex:exittime}
The function $k\mapsto k^2$ has the right second difference. Apply optional reasoning by conditioning on the first step, as for $e_k$ in the ruin section.
\end{hint}
\begin{hint}{ex:electric}
The current out of $a$ along the edges to its neighbours is $\sum_{y\sim a}h(y)$. The probability of hitting $b$ before returning is that current divided by the degree, because the first step chooses a neighbour uniformly and then hits with probability $h$.
\end{hint}

\section{Solutions}

\begin{solution}{ex:ballotwalk}
\idea The ballot theorem counts paths that stay positive. Reflection compares paths that touch a barrier with unrestricted paths.
\writeup Theorem~\ref{thm:ballot} counts sequences that stay strictly positive and end at $a-b>0$. Theorem~\ref{thm:rwreflect} counts sequences that touch a level. They are converses in the following sense. A path from $0$ to $k>0$ in $n$ steps stays at least $1$ after time $0$ precisely when its first step is $+1$ and the shifted path does not touch $0$. The number of such paths was computed in the proof of Theorem~\ref{thm:ballot} by subtracting the reflected paths, which is the same subtraction as in Theorem~\ref{thm:rwreflect} with barrier $0$ after the first step. The hypothesis $a>b$ in the ballot theorem is the statement that the end $k=a-b$ is positive, so a path can stay positive and still finish there.
\end{solution}

\begin{solution}{ex:maxlaw}
\idea Sum the pointwise identity over the admissible ends.
\writeup For each $b\le a$ of the right parity, Theorem~\ref{thm:rwreflect} gives $P(M_n\ge a,\,S_n=b)=P(S_n=2a-b)$. Sum over those $b$. As $b$ runs through the admissible integers $\le a$, the point $2a-b$ runs through the admissible integers $\ge a$. Each such integer arises once. The term $b=a$ produces $P(S_n=a)$, and the terms $b<a$ produce $P(S_n>a)$. Hence $P(M_n\ge a)=P(S_n\ge a)+P(S_n>a)$.
\end{solution}

\begin{solution}{ex:timepositive}
\idea A first return is a path to zero whose interior does not touch zero. Reflection removes the paths that do touch.
\writeup Any path with $S_{2m}=0$ has a first return to zero at some time $2j\le 2m$. The number of paths from $0$ to $0$ in $2m$ steps whose first return is at $2m$ equals the number of paths that stay strictly positive on $\{1,\ldots,2m-1\}$ and end at $0$, times $2$ for the sign, or equivalently the paths from $1$ to $1$ in $2m-2$ steps that stay positive. The reflection argument of Theorem~\ref{thm:ballot} gives the count $2^{2m}f_{2m}$ with $f_{2m}=u_{2m-2}/(2m-1)$. The standard identity $f_{2m}=u_{2m-2}-u_{2m}$ is the same number, because $u_{2m}=\binom{2m}{m}2^{-2m}$ and $u_{2m-2}-u_{2m}=u_{2m}/(2m-1)$ follows from $\binom{2m}{m}/\binom{2m-2}{m-1}=(4m-2)/m$, a short algebraic check. Thus $P(\text{first return at }2m)=u_{2m}/(2m-1)$.
\end{solution}

\begin{solution}{ex:unfair}
\idea Both the formula and three linear equations describe the same function.
\writeup Here $q/p=1/2$. Then $h_k=((1/2)^k-1)/((1/2)^3-1)=(1-2^{-k})/(1-1/8)=(1-2^{-k})/(7/8)$. So $h_1=(1/2)/(7/8)=4/7$ and $h_2=(3/4)/(7/8)=6/7$. The system is $h_1=(2/3)h_2$, $h_2=(2/3)\cdot1+(1/3)h_1$, with $h_0=0$ and $h_3=1$. Substitution gives the same values.
\end{solution}

\begin{solution}{ex:meet}
\idea The difference of the two positions is a symmetric walk on the even integers, started at $2$, and that walk is recurrent.
\writeup Let the walks be $S$ and $2+S'$. Their difference at time $n$ is $S_n-S'_n-2$. Each step of the difference is $0$ with probability $1/2$ and $\pm2$ with probability $1/4$. The embedded walk that ignores the stays is a fair simple random walk on $\Z$, scaled by $2$, and it visits every site infinitely often by Theorem~\ref{thm:polya}. The stays do not prevent visits. Starting at $-2$ in the difference, the walk reaches $0$ almost surely.
\end{solution}

\begin{solution}{ex:escape}
\idea The first-return probabilities decay like $m^{-3/2}$, and $m$ times that is like $m^{-1/2}$.
\writeup By the previous exercise, $P(\tau=2m)=u_{2m}/(2m-1)\sim (1/\sqrt{\pi m})/(2m)=O(m^{-3/2})$. Then $2m\,P(\tau=2m)\sim 1/\sqrt{\pi m}$, and $\sum m^{-1/2}$ diverges. The expectation of $\tau$ is therefore infinite.
\end{solution}

\begin{solution}{ex:exittime}
\idea The function $k^2$ is harmonic up to a constant drift of $1$ per step, on the interval $(-R,R)$.
\writeup Let $e_k=\mathbb E_k[\tau]$ for $|k|\le R$, with $e_R=e_{-R}=0$. Conditioning gives $e_k=1+(e_{k+1}+e_{k-1})/2$ for $|k|<R$. The function $R^2-k^2$ satisfies the same equation, because its second difference is $-2$, and it vanishes at $\pm R$. Uniqueness follows as in Proposition~\ref{prop:harmonic}: the difference of two solutions would be harmonic and zero on the boundary. Thus $e_0=R^2$.
\end{solution}

\begin{solution}{ex:electric}
\idea Harmonic functions are the voltages of unit resistors. The escape probability is the current out of $a$ divided by the degree.
\writeup (a) This is Proposition~\ref{prop:harmonic}.
(b) Assign voltage $h(x)$ to vertex $x$. The current from $x$ to a neighbour $y$ is $h(x)-h(y)$ if each edge has resistance $1$. At $x\notin\{a,b\}$ the sum of currents is $\deg(x)h(x)-\sum_{y\sim x}h(y)=0$, because $h$ is harmonic. No current is created there. The boundary values are the imposed voltages.
(c) The total current leaving $b$ and entering the network equals the current into $a$, which is $\sum_{y\sim a}(h(y)-h(a))=\sum_{y\sim a}h(y)$. Call it $I$, and set $R=1/I$ when $I>0$, which it is because $h$ is not identically zero on a connected graph. Starting at $a$, the first step moves to a uniform neighbour $y$, and from there the probability of hitting $b$ before $a$ is $h(y)$. The probability of hitting $b$ before returning to $a$ is therefore $I/\deg(a)=1/(\deg(a)R)$.
(d) On a path of length $m$ the voltages are equally spaced, $h=j/m$ at distance $j$ from $a$. The current is $1/m$, so $R=m$. Resistors in series add. On a cycle with paths of lengths $m$ and $k$, the two paths are in parallel. Each has resistance equal to its length, and the parallel law gives $1/R=1/m+1/k$, so $R=mk/(m+k)$.
\end{solution}

\section*{Sources and further reading}

\begin{enumerate}
\item G.~P\'olya, \"Uber eine Aufgabe der Wahrscheinlichkeitsrechnung betreffend die Irrfahrt im Stra\ss ennetz, \emph{Math.\ Ann.} \textbf{84} (1921), 149--160. Theorem~\ref{thm:polya}.
\item W.~Feller, \emph{An Introduction to Probability Theory and Its Applications}, vol.~I, 3rd ed., Wiley, 1968, Chapter~III. Reflection and ruin.
\end{enumerate}
